\chapter{What this manual is, and what it has proven}\label{ch:intro}

This manual is for the person who has to change the ATE application: a Python developer who may not
have written a Test Schedule, and who may never have threaded a Qt application before. The rules here
are not style rules. Each one exists because a run broke in a specific way, and the fix is a rule
about where code may go and what it may call. A rule you can recite but not explain is a rule the next
change breaks.

Two files make this manual, and they say the same things:

\begin{itemize}
  \item this PDF, built from \texttt{docs/manuals/developer/} by \texttt{python docs/manuals/build.py developer};
  \item \texttt{docs/manuals/llm/developer.md}, written for an agent rather than converted from this
    text: it is shorter, it is in the order an agent works, and the checks below apply to it too.
\end{itemize}

\section{What is checked, and how}

\texttt{tests/test_manuals.py} holds the manuals to the code. It does not read them for style; it
derives lists of names from the code and fails when a manual stops naming one:

\begin{itemize}
  \item every module under \texttt{src/ate/}, derived with \texttt{pkgutil} from the package itself;
  \item every Qt signal declared on \texttt{RunEngine} or \texttt{HeadlessRunSession}, read off the
    classes;
  \item every table in \texttt{ate.storage.schema}, read from the SQLAlchemy \texttt{metadata};
  \item every Device driver class and the Device Type name it serves, taken from the classes that
    subclass \texttt{ate.devices.registry.Driver}.
\end{itemize}

The same test extracts every \texttt{src/...} path and every dotted name in backticks from the two
developer manuals and looks them up: a file that does not exist, or a module that does not have the
attribute named, fails the suite. Commands in a shell listing are each either run by the test suite in
a temporary Station or listed with the reason they cannot be run here. The two worked examples,
adding a Device type (chapter~\ref{ch:device-type}) and adding a schema version
(chapter~\ref{ch:schema-version}), are executed against a temporary copy of the source: the steps the
manual gives are the steps the test performs.

\section{How to run things}

Nothing here needs a display. Every Qt test runs with the offscreen platform:

\begin{lstlisting}[language=ATEShell]
QT_QPA_PLATFORM=offscreen python -m pytest -q -p no:cacheprovider
\end{lstlisting}

That is the command CI runs, and the one this manual was written against. Two gates are off by
default because they take minutes or need Docker:

\begin{lstlisting}[language=ATEShell]
ATE_BUILD_TESTS=1 pytest -m build      # builds the real PyInstaller bundle
ATE_DOCKER_TESTS=1 pytest -m build     # runs the debian:stable container smokes
\end{lstlisting}

Both are described in chapter~\ref{ch:tests}, and both were run for this manual.

\section{What is proven here, and what is not}

Proven on this machine, by the tests named in the chapters that state them:

\begin{itemize}
  \item the whole suite on Linux with the offscreen platform, repeated until it was green;
  \item the threading behaviour of every signal the engine and the session emit, and the thread each
    slot runs on (\texttt{tests/test_signal_threads_32a.py});
  \item that an abort ends a Device call on every Device Type, and how long it takes
    (\texttt{tests/test_device_abort_32b.py});
  \item that a new Device driver and a new schema version work when added the way this manual says.
\end{itemize}

Not proven, and this manual says so wherever it matters:

\begin{itemize}
  \item \textbf{real hardware.} Every Device test here talks to a loopback peer in
    \texttt{tests/_fixtures.py} or a pseudo-terminal. A GPIB instrument, a real serial port and a
    real VISA stack have never been touched by the suite.
  \item \textbf{real Windows sessions.} Windows behaviour is asserted by the same tests, but only
    GitHub Actions has ever run them on Windows, and only as CI runs.
  \item \textbf{the cause of the 32b stall.} One Windows CI run stalled a Test Run for tens of
    seconds. The fixes are proven; the cause is a belief about the disk, not a measurement.
    Chapter~\ref{ch:failed-run} says what a future failure would show.
\end{itemize}

\cando{tell which claim in this manual rests on a test you can run, and which rests on a machine that
is not here.}
